\kchapter{Before Kubernetes: Applications and Containers}{docker}

\begin{objectives}
\begin{itemize}
  \item why ``it works on my machine'' is the oldest problem in software;
  \item what a container is, and how it differs from a virtual machine;
  \item the vocabulary: image, container, registry, tag, layer;
  \item how to write a Dockerfile and build, run, and push an image.
\end{itemize}
\end{objectives}

Kubernetes is a system for running \emph{containers}. Before we can understand
Kubernetes, we need to understand what it is managing. This chapter is your
crash course.

\section{The problem: ``it works on my machine''}

Imagine you have written a small web API in Python. On your laptop it runs
perfectly. You copy it to a server and it crashes: the server has a different
Python version, a missing system library, and a different configuration file.
A teammate tries it on a Mac and gets a third set of errors.

Every application depends on much more than its own source code:

\begin{itemize}
  \item a language runtime (Python 3.12, Node 22, Java 21, \ldots);
  \item libraries and their exact versions;
  \item operating-system packages (OpenSSL, libpq, fonts, timezone data);
  \item environment variables and configuration files;
  \item a particular directory layout and file permissions.
\end{itemize}

Historically, teams solved this with long installation documents and careful
manual server setup. That is slow, fragile, and impossible to repeat exactly.

\section{Virtual machines: the first answer}

A \kterm{virtual machine} (VM) emulates a whole computer. You install a full
operating system inside it, then your application. A VM image can be copied
anywhere and behaves the same. The cost: every VM carries a complete operating
system kernel, takes gigabytes of disk, and needs a minute or more to boot.

\section{Containers: a lighter answer}

A \kterm{container} packages an application \emph{together with everything it
needs to run} --- runtime, libraries, OS packages, config --- but it does
\emph{not} include its own kernel. All containers on a machine share the host's
Linux kernel. The kernel uses two features to keep containers apart:

\begin{itemize}
  \item \kterm{namespaces} give each container its own private view of
        processes, network interfaces, mount points, and hostnames;
  \item \kterm{cgroups} (control groups) limit how much CPU, memory, and I/O a
        container may use.
\end{itemize}

So a container is really just an ordinary Linux process that has been given
blinkers (namespaces) and a budget (cgroups).

\begin{center}
\resizebox{0.92\linewidth}{!}{%
\begin{tikzpicture}[
  layer/.style={draw=none, fill=#1, text=white, minimum height=0.8cm,
                font=\small\bfseries, rounded corners=2pt},
  app/.style={draw=none, fill=#1, text=white, minimum height=0.8cm,
              minimum width=1.9cm, font=\footnotesize\bfseries, rounded corners=2pt}]
  % VMs
  \node[font=\bfseries, text=k8snavy] at (3.1,5.3) {Virtual Machines};
  \node[layer=k8sgray, minimum width=6.4cm] at (3.1,0) {Hardware};
  \node[layer=k8sgray!80, minimum width=6.4cm] at (3.1,0.9) {Hypervisor};
  \foreach \x/\n in {1.0/A, 3.1/B, 5.2/C}{
    \node[app=k8sdeep] at (\x,1.8) {Guest OS};
    \node[app=k8smid, text=k8snavy] at (\x,2.7) {Libs};
    \node[app=k8sblue] at (\x,3.6) {App \n};
  }
  \draw[decorate, decoration={brace, amplitude=6pt, mirror}, k8sgray]
       (6.45,1.4) -- (6.45,4.0) node[midway, right=8pt, font=\scriptsize, text=k8sgray, align=left]{GBs each,\\minutes to boot};
  % Containers
  \node[font=\bfseries, text=k8snavy] at (11.6,5.3) {Containers};
  \node[layer=k8sgray, minimum width=6.4cm] at (11.6,0) {Hardware};
  \node[layer=k8sgray!80, minimum width=6.4cm] at (11.6,0.9) {Host OS + shared Linux kernel};
  \node[layer=k8sdeep, minimum width=6.4cm] at (11.6,1.8) {Container runtime (containerd)};
  \foreach \x/\n in {9.5/A, 11.6/B, 13.7/C}{
    \node[app=k8smid, text=k8snavy] at (\x,2.7) {Libs};
    \node[app=k8sblue] at (\x,3.6) {App \n};
  }
  \draw[decorate, decoration={brace, amplitude=6pt, mirror}, k8sgray]
       (14.95,2.3) -- (14.95,4.0) node[midway, right=8pt, font=\scriptsize, text=k8sgray, align=left]{MBs each,\\start in ms};
\end{tikzpicture}}
\end{center}

\begin{analogy}
A virtual machine is a \emph{detached house}: its own foundations, plumbing and
roof. A container is an \emph{apartment} in a building: it has its own locked
front door and furniture, but shares the foundations and plumbing (the kernel)
with its neighbours. Apartments are cheaper and quicker to build.
\end{analogy}

\section{The container vocabulary}

\begin{center}
\ktablestart
\begin{tabularx}{\linewidth}{P{2.8cm} L}
\kheader \kth{Term} & \kth{Meaning} \\
Image & A read-only, versioned package of an application and its
        filesystem. Think of it as a \emph{class} or a \emph{recipe}. \\
Container & A running instance of an image. Think \emph{object} or
            \emph{the cooked meal}. Many containers can run from one image. \\
Layer & Images are built in stacked, cached layers. Changing one line of your
        Dockerfile only rebuilds the layers after it. \\
Registry & A server that stores images: Docker Hub, GitHub Container Registry
           (\kc{ghcr.io}), Amazon ECR, Google Artifact Registry, Harbor\ldots \\
Tag & A human label for an image version, e.g.\ \kc{nginx:1.27}. \\
Digest & An immutable content hash, e.g.\ \kc{nginx@sha256:4c0f\ldots}. A
         tag can be moved; a digest cannot. \\
Dockerfile & The text recipe used to build an image. \\
OCI & The Open Container Initiative: the open standard for image and runtime
      formats. Docker, Podman, containerd, and Kubernetes all speak OCI. \\
\end{tabularx}
\end{center}

\section{Your first Dockerfile}

Let us containerise a tiny Python web server. Create a folder with two files.

\begin{pycode}{app.py}
from http.server import BaseHTTPRequestHandler, HTTPServer
import os, socket

class Hello(BaseHTTPRequestHandler):
    def do_GET(self):
        msg = f"Hello from {socket.gethostname()} (version {os.getenv('VERSION', '1')})\n"
        self.send_response(200)
        self.end_headers()
        self.wfile.write(msg.encode())

HTTPServer(("0.0.0.0", 8080), Hello).serve_forever()
\end{pycode}

\begin{dockercode}{Dockerfile}
# 1. Start from an official, small base image
FROM python:3.12-slim

# 2. Work inside /app in the image
WORKDIR /app

# 3. Copy our code into the image
COPY app.py .

# 4. Run as a non-root user (good security habit)
RUN useradd --uid 10001 appuser
USER 10001

# 5. Document the port and define the start command
EXPOSE 8080
CMD ["python", "app.py"]
\end{dockercode}

Now build, run, and test it:

\begin{bashcode}{terminal}
# Build an image called hello with tag 1.0 from the current folder (.)
docker build -t hello:1.0 .

# Run it in the background (-d), mapping host port 8080 to container port 8080
docker run -d --name hello -p 8080:8080 hello:1.0

curl http://localhost:8080
# Hello from 3f2a9c1b7e44 (version 1)

docker ps            # list running containers
docker logs hello    # see its output
docker stop hello && docker rm hello
\end{bashcode}

\subsection{Pushing to a registry}

Kubernetes nodes pull images from a registry, so images must be pushed
somewhere the cluster can reach.

\begin{bashcode}{terminal}
docker login ghcr.io -u YOUR_GITHUB_USER
docker tag hello:1.0 ghcr.io/YOUR_GITHUB_USER/hello:1.0
docker push ghcr.io/YOUR_GITHUB_USER/hello:1.0
\end{bashcode}

\section{Good image habits}

\begin{itemize}
  \item \textbf{Use small, official base images} (\kc{-slim}, \kc{-alpine},
        or ``distroless''). Less software means fewer vulnerabilities.
  \item \textbf{Pin versions.} \kc{python:3.12-slim} is better than
        \kc{python:latest}; a digest is better still.
  \item \textbf{Order layers from least to most frequently changed}: install
        dependencies first, copy source code last, so rebuilds are fast.
  \item \textbf{Use multi-stage builds} to compile in one stage and copy only
        the result into a small runtime stage.
  \item \textbf{Run as non-root.} Kubernetes can enforce this for you later.
  \item \textbf{One process per container.} Let the platform handle restarts
        and scaling instead of running a process supervisor inside.
  \item \textbf{Log to stdout/stderr}, never to files inside the container.
\end{itemize}

\section{Why containers are not enough}

Containers solve packaging. But once you run real systems, new questions appear:

\begin{itemize}
  \item You have 40 containers across 10 servers. Which server should a new one go on?
  \item A server dies at 3 a.m. Who restarts its containers somewhere else?
  \item Traffic triples on Black Friday. Who starts more copies, and who
        removes them afterwards?
  \item How do containers find each other when their IP addresses keep changing?
  \item How do you roll out version 2 without downtime, and roll back if it is broken?
  \item Where do passwords and configuration live?
\end{itemize}

Answering these questions automatically is called \kterm{container
orchestration}. And the de-facto standard orchestrator is Kubernetes.

\begin{recap}
\begin{itemize}
  \item Containers package an app with its dependencies but share the host kernel.
  \item An \emph{image} is the recipe; a \emph{container} is a running instance.
  \item Images are built from a Dockerfile and stored in a registry.
  \item Running many containers reliably across many machines needs an
        orchestrator --- that is Kubernetes.
\end{itemize}
\end{recap}
